The scarcity of high-quality, logically annotated video datasets remains a primary bottleneck in advancing Multi-Modal Large Language Models (MLLMs) for the medical domain. Traditional manual annotation is prohibitively expensive and non-scalable, while existing synthetic methods often suffer from stochastic hallucinations and a lack of logical interpretability. To address these challenges, we introduce \textbf{\PipelineName}, a novel neuro-symbolic data engineering framework that formalizes benchmark synthesis as a deterministic graph traversal process.
% 高质量、逻辑标注的视频数据集的匮乏仍然是推进医学领域多模态大型语言模型（MLLM）发展的主要瓶颈。传统的人工标注成本高昂且难以扩展，而现有的合成方法往往存在随机幻觉和逻辑可解释性不足的问题。为了应对这些挑战，我们提出了\PipelineName，这是一个新颖的神经符号数据工程框架，它将基准合成形式化为一个确定性的图遍历过程。
Unlike black-box generative approaches, \PipelineName~ extracts structured visual primitives (e.g., surgical instruments, anatomical boundaries) from raw video streams and instantiates them into a dynamic Spatiotemporal Knowledge Graph. By anchoring query generation to valid paths within this graph, we enforce a rigorous Chain-of-Thought (CoT) provenance for every synthesized benchmark item. We instantiate this pipeline to produce \BenchmarkName, a large-scale medical video reasoning benchmark exhibiting fine-grained temporal selectivity and multi-hop logical complexity.
% 与黑盒生成方法不同，\PipelineName 从原始视频流中提取结构化的视觉基元（例如，手术器械、解剖边界），并将其实例化为动态的时空知识图谱。通过将查询生成锚定到该图谱中的有效路径，我们为每个合成的基准测试项强制执行严格的“思维链”（CoT）溯源。我们实例化此流程生成 \BenchmarkName，这是一个大规模的医学视频推理基准测试，展现出细粒度的时间选择性和多跳逻辑复杂性。
Comprehensive evaluations demonstrate that our automated pipeline generates query workloads with complexity comparable to expert-curated datasets. Furthermore, a logic alignment analysis reveals a high correlation between the prescribed graph topology and the reasoning steps of state-of-the-art MLLMs, validating the system's capability to encode verifiable logic into visual-linguistic benchmarks. This work paves the way for scalable, low-cost construction of robust evaluation protocols in critical domains.
% 全面的评估表明，我们的自动化流程生成的查询工作负载的复杂度与专家精心整理的数据集相当。此外，逻辑对齐分析揭示了预设图拓扑结构与最先进的多层逻辑模型（MLLM）推理步骤之间的高度相关性，验证了系统将可验证逻辑编码到视觉语言基准测试中的能力。这项工作为在关键领域构建可扩展、低成本的稳健评估协议铺平了道路。